\begin{center}
\begin{tikzpicture}[x=.88cm,y=.88cm,font=\small]
  \fill[paper] (0,0) rectangle (10.2,5.2);
  \draw[blue!55,line width=1pt] (.3,4.3) .. controls (2.6,4.1) and (5.2,4.6) .. (9.8,4.1);
  \node[text=blue!70!black,above] at (5.0,4.25) {黄河方向};
  \fill[dynasty!23] (3.8,2.2) -- (5.9,2.1) -- (6.1,3.8) -- (4.0,4.0) -- cycle;
  \node[align=center] at (5.0,3.0) {周王畿\\洛邑方向};
  \fill[green!20] (6.8,1.0) -- (9.3,1.1) -- (9.5,3.3) -- (7.2,3.5) -- cycle;
  \node[align=center] at (8.2,2.2) {郑\\中原要冲};
  \fill[timeline] (6.4,3.1) circle (.14);
  \node[above,align=center] at (6.4,3.2) {繻葛\\战场方向};
  \draw[-{Stealth[length=2mm]},ultra thick,color=dynasty] (4.9,2.7) -- (6.25,3.05);
  \draw[-{Stealth[length=2mm]},thick,color=green!60!black] (8.0,2.55) -- (6.55,3.05);
  \node[text=dynasty,below] at (5.5,2.55) {王师东出};
  \node[text=green!60!black,below] at (7.3,2.6) {郑军应战};
\end{tikzpicture}

\smallskip
\small\textit{图  繻葛之战空间示意：突出周王畿、郑国与战场方向的关系；不表示精确行军路线或疆界。}
\end{center}
